\documentclass[12pt]{article}
\usepackage[T1]{fontenc}
\usepackage[utf8]{inputenc}
\usepackage{lmodern}
\usepackage{textcomp}
\usepackage{enumitem}
\usepackage{geometry}
\geometry{margin=1in}
\begin{document}
\begin{center}
Answer Key - Biology - Graduate Level Assessment 16 \
\small (Answers are ordered exactly as in the question paper.)
\end{center}

\section*{Multiple Choice}
\begin{enumerate}[label=\arabic*.]
\item \textbf{Answer:} D
\\ \textit{Options:} A) a cell wall made of cellulose; B) lysosomes; C) large vacuoles that store water; D) centrioles within centrosomes
\\ \textit{Note:} Animal cells contain centrioles organized within centrosomes, which are crucial for cell division and are absent in plant cells, highlighting a key structural difference between the two cell types.
\item \textbf{Answer:} E
\\ \textit{Options:} A) auxins; B) neurotransmitters; C) enzymes; D) hormones; E) pheromones
\\ \textit{Note:} Pheromones are chemical substances produced by organisms that trigger specific physiological or behavioral responses in other individuals of the same species, thereby playing a crucial role in communication and social interaction among them.
\item \textbf{Answer:} A
\\ \textit{Options:} A) Carbon dioxide bubbles forming in the blood causing acidosis; B) Rapid increase in barometric pressure causing the bends; C) Formation of helium bubbles in the bloodstream due to rapid ascent; D) Increased salinity in the water leading to dehydration and joint pain; E) Oxygen bubbles forming in the blood leading to hypoxia
\\ \textit{Note:} The bends, or decompression sickness, occurs when a diver ascends too quickly, causing dissolved gases, primarily nitrogen, to form bubbles in the bloodstream, leading to tissue damage and acidosis due to the release of carbon dioxide from the tissues.
\item \textbf{Answer:} A
\\ \textit{Options:} A) homeotic genes; B) cell division; C) signal transduction; D) cell differentiation; E) apoptosis
\\ \textit{Note:} Homeotic genes are primarily responsible for determining the identity of body parts and structures during development, rather than directly influencing the processes involved in morphogenesis, which include cell shape changes and tissue organization.
\item \textbf{Answer:} A
\\ \textit{Options:} A) multiple copies of a single chromosome; B) not visible even under a light microscope; C) produced only during mitosis; D) produced by repeated rounds of nuclear division without DNA replication; E) single copies of multiple chromosomes
\\ \textit{Note:} The correct answer is due to the fact that polytene chromosomes are formed by the replication of a single chromosome, resulting in multiple identical copies that remain aligned in parallel, which can be visualized under a light microscope.
\end{enumerate}

\section*{True / False}
\begin{enumerate}[label=\arabic*.]
\setcounter{enumi}{5}
\item \textbf{Answer:} False
\\ \textit{Options:} A) True; B) False
\\ \textit{Note:} While more than a quarter of the middle-aged adult population said they were interested in PAS insurance, actual participation would be highly dependent on premium rates. The current lack of publicly subsidized insurance for long-term care and personal assistance services remains a serious gap in the disability service system.
\item \textbf{Answer:} True
\\ \textit{Options:} A) True; B) False
\\ \textit{Note:} The gluten-free diet has minimal deficiencies, similar to those present in the diet with gluten, with an improvement in the lipid profile by increasing the proportion of monounsaturated fatty acids to the detriment of saturated fatty acids.
\item \textbf{Answer:} True
\\ \textit{Options:} A) True; B) False
\\ \textit{Note:} Ketamine sedation was successful and well tolerated in all cases. The use of atropine as an adjunct for intramuscular ketamine sedation in children significantly reduces hypersalivation and may lower the incidence of post-procedural vomiting. Atropine is associated with a higher incidence of a transient rash. No serious adverse events were noted.
\item \textbf{Answer:} False
\\ \textit{Options:} A) True; B) False
\\ \textit{Note:} While FS was commonly used in patients undergoing thyroidectomy at our institution, in no patient over the last decade did FS correctly alter the intraoperative management. Given the time required to perform FS and the cost associated with it, we believe that routine FS should not be performed in these patients.
\item \textbf{Answer:} True
\\ \textit{Options:} A) True; B) False
\\ \textit{Note:} It is important to ensure that new mothers are adequately informed about topics important to them while in hospital. The findings highlight the need for accessible and appropriate community-based information resources for women in the postpartum period, especially for those of low socioeconomic status.
\end{enumerate}

\section*{Long Form Response}
\begin{enumerate}[label=\arabic*.]
\setcounter{enumi}{10}
\item \textbf{Answer:} The term 'gymnosperm' refers to plants that produce 'naked seeds,' which are not enclosed by an ovary or fruit, while 'angiosperm' denotes plants with 'enclosed seeds' that develop within a fruit. Conifers, a prominent group of gymnosperms, are characterized by their adaptations to colder climates, which include needle-like leaves that minimize water loss and a reproductive strategy relying on wind pollination and the formation of cone-like structures for seed dispersal. In contrast, flowering plants exhibit a more complex reproductive strategy that involves the development of flowers to attract pollinators, the mechanism of fertilization leading to seed formation within fruits, and a wider variety of dispersal methods. These distinctions highlight how conifers and angiosperms have evolved different adaptations and reproductive strategies to thrive in their respective environments.
\\ \textit{Note:} The answer correctly distinguishes gymnosperms and angiosperms based on their seed characteristics and elaborates on the adaptations and reproductive strategies of conifers compared to flowering plants, highlighting the ecological significance of these traits.
\item \textbf{Answer:} The cell cycle consists of several distinct phases, each with specific functions crucial for proper cell division and growth. The M phase, or mitosis, is where active cell division occurs, resulting in the formation of two daughter cells. Before this, cells may enter G₀, a phase of temporary arrest where they can remain metabolically active but not proliferative, allowing for cellular differentiation or a pause in the cycle. In the G₁ phase, cells engage in growth and organelle duplication, ensuring that the necessary cellular machinery is present for subsequent division. The S phase is critical for DNA synthesis, during which the cell replicates its genetic material and also involves extensive RNA transcription to prepare for protein synthesis. Finally, in the G₂ phase, the cell undergoes further growth and prepares for mitosis by ensuring all DNA is intact and that the cell has adequate resources for division, thus emphasizing the intricate coordination required for successful cell cycle progression.
\\ \textit{Note:} The answer correctly identifies and elaborates on the key phases of the cell cycle-M, G₀, G₁, S, and G₂-each of which plays a vital role in regulating cell growth, organelle duplication, and the preparation necessary for successful cell division.
\end{enumerate}

\vfill
\noindent\textit{End of Answer Key}
\end{document}